\subsection{Refinement and coarse-graining channels}

% **Local fix (zero-weight quotient):** the source quotient is completed on
% zero-weight pairs. Documented in
% docs/paper-gaps/cpgsv17_mpdo_zero_weight_preparation_completion.tex.
\begin{theorem}[Positivity and zero-weight neighboring operators]
    \label{thm:mpdo_physical_sector_three_site_neighboring_operator_properties}
    \lean{MPOTensor.PhysicalSectorFactorization.NeighboringTraceFactorization.sector_nonempty,
      MPOTensor.PhysicalSectorFactorization.NeighboringTraceFactorization.threeSiteMiddleIndex_nonempty}
    \lean{MPOTensor.PhysicalSectorFactorization.NeighboringTraceFactorization.threeSiteNeighboringOperator_pos,
      MPOTensor.PhysicalSectorFactorization.NeighboringTraceFactorization.threeSiteNeighboringOperator_eq_zero_of_mul_eq_zero}
    \leanok
    \uses{def:mpdo_physical_sector_three_site_neighboring_operator,
      def:mpdo_neighboring_trace_factorization,
      thm:mpdo_physical_sector_three_site_neighboring_operator_trace}
    The sector set and every middle-subspin space are nonempty. Each
    $\Omega_{k,h}$ is positive semidefinite, and
    \[
      a_kb_h=0 \quad\Longrightarrow\quad \Omega_{k,h}=0.
    \]
\end{theorem}

\begin{definition}[The three-site neighboring-state preparation $\mathcal T_1$]
    \label{def:mpdo_physical_sector_t1_preparation}
    \lean{MPOTensor.PhysicalSectorFactorization.NeighboringTraceFactorization.normalizedThreeSiteNeighboringDensity}
    \lean{MPOTensor.PhysicalSectorFactorization.NeighboringTraceFactorization.inactiveThreeSiteNeighboringDensity}
    \lean{MPOTensor.PhysicalSectorFactorization.NeighboringTraceFactorization.completedThreeSiteNeighboringDensity}
    \lean{MPOTensor.PhysicalSectorFactorization.NeighboringTraceFactorization.completedThreeSitePreparationMap}
    \lean{MPOTensor.PhysicalSectorFactorization.NeighboringTraceFactorization.completedThreeSiteControlledMap}
    \lean{MPOTensor.PhysicalSectorFactorization.NeighboringTraceFactorization.t1Map}
    \leanok
    \uses{def:mpdo_physical_sector_three_site_neighboring_operator,
      def:mpdo_neighboring_trace_factorization,
      thm:mpdo_physical_sector_three_site_neighboring_operator_properties}
    On a sector pair with $a_kb_h\ne0$, set
    \[
      \widehat\Omega_{k,h}=(a_kb_h)^{-1}\Omega_{k,h}.
    \]
    If $a_kb_h=0$, positivity and vanishing trace imply
    $\Omega_{k,h}=0$; choose any density operator on that summand and denote
    the resulting completed density by $\overline\Omega_{k,h}$. The sectorwise
    preparation is
    \[
      X\longmapsto X\otimes\overline\Omega_{k,h}.
    \]
    Taking the orthogonal direct sum over the outer sectors $(k,h)$ gives the
    preparation stage $\mathcal T_1$ of
    \cite[Appendix~C.2, lines~1523--1535]{Cirac2016MPDO_arXiv}.
    The zero-weight choice makes the source map total without imposing a
    nonzero-weight hypothesis.
\end{definition}

\begin{theorem}[Properties of the three-site neighboring densities]
    \label{thm:mpdo_physical_sector_t1_density_properties}
    \lean{MPOTensor.PhysicalSectorFactorization.NeighboringTraceFactorization.normalizedThreeSiteNeighboringDensity_pos,
      MPOTensor.PhysicalSectorFactorization.NeighboringTraceFactorization.trace_normalizedThreeSiteNeighboringDensity}
    \lean{MPOTensor.PhysicalSectorFactorization.NeighboringTraceFactorization.inactiveThreeSiteNeighboringDensity_posDef,
      MPOTensor.PhysicalSectorFactorization.NeighboringTraceFactorization.inactiveThreeSiteNeighboringDensity_trace}
    \lean{MPOTensor.PhysicalSectorFactorization.NeighboringTraceFactorization.completedThreeSiteNeighboringDensity_pos,
      MPOTensor.PhysicalSectorFactorization.NeighboringTraceFactorization.trace_completedThreeSiteNeighboringDensity,
      MPOTensor.PhysicalSectorFactorization.NeighboringTraceFactorization.completedThreeSiteNeighboringDensity_eq_normalized}
    \lean{MPOTensor.PhysicalSectorFactorization.NeighboringTraceFactorization.completedThreeSitePreparationMap_eq_normalized}
    \lean{MPOTensor.PhysicalSectorFactorization.NeighboringTraceFactorization.completedThreeSiteControlledMap_sameBlock_apply}
    \leanok
    \uses{def:mpdo_physical_sector_t1_preparation}
    The normalized density is positive semidefinite and has trace one on
    every active pair. The density chosen on an inactive pair is positive
    definite and has trace one. Hence $\overline\Omega_{k,h}$ is positive
    semidefinite with trace one for every pair, and agrees with
    $\widehat\Omega_{k,h}$ on active pairs. The corresponding sectorwise and
    controlled maps prepare these densities on their diagonal summands.
\end{theorem}

\begin{theorem}[$\mathcal T_1$ is trace-preserving completely positive]
    \label{thm:mpdo_physical_sector_t1_preparation_cptp}
    \lean{MPOTensor.PhysicalSectorFactorization.NeighboringTraceFactorization.completedThreeSitePreparationMap_isKrausCPTP}
    \lean{MPOTensor.PhysicalSectorFactorization.NeighboringTraceFactorization.completedThreeSiteControlledMap_isKrausCPTP}
    \lean{MPOTensor.PhysicalSectorFactorization.NeighboringTraceFactorization.t1Map_isKrausCPTP}
    \leanok
    \uses{def:mpdo_physical_sector_t1_preparation,
      thm:mpdo_physical_sector_t1_density_properties}
    Each completed preparation is trace-preserving and completely positive.
    The same holds for their orthogonally controlled direct sum.
\end{theorem}

\begin{proof}
    \leanok
    On an active pair,
    $\tr(\widehat\Omega_{k,h})=(a_kb_h)^{-1}a_kb_h=1$; the chosen density on
    a zero-weight pair also has trace one. Preparation by a positive operator
    of trace one is trace-preserving and completely positive. The orthogonal
    control discards coherences between distinct outer-sector pairs and
    applies these preparations on the diagonal summands.
\end{proof}

\begin{definition}[The refinement channel $\mathcal T$]
    \label{def:mpdo_physical_sector_t_map}
    \lean{MPOTensor.PhysicalSectorFactorization.NeighboringTraceFactorization.tMap}
    \leanok
    \uses{def:mpdo_t0_map,
      def:mpdo_physical_sector_t1_preparation,
      def:mpdo_t2_map}
    Define the channel from two physical sites to three physical sites by
    \[
      \mathcal T=\mathcal T_2\mathcal T_1\mathcal T_0.
    \]
    Thus $\mathcal T_0$ traces the neighboring subspins, $\mathcal T_1$
    adjoins the completed three-site neighboring density, and $\mathcal T_2$
    restores the three complete physical sectors. This is the refinement map
    in \cite[Appendix~C.2, lines~1522--1545]{Cirac2016MPDO_arXiv}.
\end{definition}

\begin{theorem}[$\mathcal T$ is trace-preserving completely positive]
    \label{thm:mpdo_physical_sector_t_map_cptp}
    \lean{MPOTensor.PhysicalSectorFactorization.NeighboringTraceFactorization.tMap_isKrausCPTP}
    \leanok
    \uses{def:mpdo_physical_sector_t_map}
    The refinement channel $\mathcal T$ is trace-preserving and completely
    positive.
\end{theorem}

\begin{proof}
    \leanok
    \uses{thm:mpdo_t0_map_cptp,
      thm:mpdo_physical_sector_t1_preparation_cptp,
      thm:mpdo_t2_map_cptp,
      lem:is_kraus_cptp_comp}
    Each constituent map is trace-preserving and completely positive, and
    this property is preserved under composition.
\end{proof}

Applied to the two-site closure, on the active $(k,h)$-block,
$\mathcal T_0$ traces $B_k^R\otimes B_h^L$ and sends
$\mathfrak R_2(\mathcal K_2(X))_{k,h}$ to
$a_kb_hB_{k,h}(X)$ on $B_k^L\otimes B_h^R$. The preparation
$\mathcal T_1$ then adjoins
\[
  \widehat\Omega_{k,h}
  =(a_kb_h)^{-1}\bigoplus_l(\eta_{k,l}\otimes\eta_{l,h}),
\]
and $\mathcal T_2$ orders the result as
\[
  (B_k^L\otimes B_k^R)\otimes(B_l^L\otimes B_l^R)
    \otimes(B_h^L\otimes B_h^R).
\]
The refinement uses the sitewise sector coordinates $\mathfrak R_2$ and
$\mathfrak R_3$
\cite[Appendix~C.2, lines~1510--1516 and~1522--1545]
  {Cirac2016MPDO_arXiv}.
For each outer pair, write
\[
  \mathcal H^{\Omega}_{k,h}
  =\bigoplus_l(B_k^R\otimes B_l^L)\otimes(B_l^R\otimes B_h^L)
\]
for the carrier Hilbert space of $\Omega_{k,h}$. The symbols
$\widehat\Omega_{k,h}$ and $\overline\Omega_{k,h}$ denote respectively its
normalized active density and its completed density on every sector pair.
With the matrix algebras introduced above, the refinement stages have type

% Formula: T_0 : A_2 -> A_partial, T_1 : A_partial -> A_Omega, and
% T_2 : A_Omega -> A_3. Ink: each blue stroke denotes a channel, its label
% names the stage, and the fixed left-to-right axis records composition.
\par\noindent\hfil
  \begin{tenkzcd}[maps, species={channel}]
    \mathcal A_2 &
    \mathcal A_{\partial} &
    \mathcal A_{\Omega} &
    \mathcal A_3
    \tnarrow[from={(1,1)}, to={(1,2)}, species=channel]{\mathcal T_0}
    \tnarrow[from={(1,2)}, to={(1,3)}, species=channel]{\mathcal T_1}
    \tnarrow[from={(1,3)}, to={(1,4)}, species=channel]{\mathcal T_2}
  \end{tenkzcd}
\hfil\par
Here $\mathcal T_1$ adjoins $\overline\Omega_{k,h}$ on each diagonal
outer-sector summand. The composite channel has type
% Formula: T = T_2 T_1 T_0 : A_2 -> A_3. Ink: the blue stroke denotes the
% composite channel T; its endpoints are matrix algebras, not matrices.
\par\noindent\hfil
  \begin{tenkzcd}[maps, species={channel}]
    \mathcal A_2 & \mathcal A_3
    \tnarrow[from={(1,1)}, to={(1,2)}, species=channel]{\mathcal T}
  \end{tenkzcd}
\hfil\par

\begin{definition}[The neighboring-state preparation $\mathcal S_1$]
    \label{def:mpdo_physical_sector_s1_map}
    \lean{MPOTensor.PhysicalSectorFactorization.NeighboringTraceFactorization.normalizedNeighboringDensity}
    \lean{MPOTensor.PhysicalSectorFactorization.NeighboringTraceFactorization.normalizedNeighboringDensity_pos,
      MPOTensor.PhysicalSectorFactorization.NeighboringTraceFactorization.trace_normalizedNeighboringDensity,
      MPOTensor.PhysicalSectorFactorization.NeighboringTraceFactorization.normalizedNeighboringPreparationMap_isKrausCPTP}
    \lean{MPOTensor.PhysicalSectorFactorization.NeighboringTraceFactorization.inactiveNeighboringDensity,
      MPOTensor.PhysicalSectorFactorization.NeighboringTraceFactorization.inactiveNeighboringDensity_posDef,
      MPOTensor.PhysicalSectorFactorization.NeighboringTraceFactorization.inactiveNeighboringDensity_trace}
    \lean{MPOTensor.PhysicalSectorFactorization.NeighboringTraceFactorization.completedNeighboringDensity}
    \lean{MPOTensor.PhysicalSectorFactorization.NeighboringTraceFactorization.completedNeighboringDensity_pos,
      MPOTensor.PhysicalSectorFactorization.NeighboringTraceFactorization.trace_completedNeighboringDensity,
      MPOTensor.PhysicalSectorFactorization.NeighboringTraceFactorization.completedNeighboringDensity_eq_normalized}
    \lean{MPOTensor.PhysicalSectorFactorization.NeighboringTraceFactorization.completedNeighboringPreparationMap,
      MPOTensor.PhysicalSectorFactorization.NeighboringTraceFactorization.completedNeighboringPreparationMap_isKrausCPTP,
      MPOTensor.PhysicalSectorFactorization.NeighboringTraceFactorization.completedNeighboringPreparationMap_eq_normalized}
    \lean{MPOTensor.PhysicalSectorFactorization.NeighboringTraceFactorization.completedNeighboringControlledMap}
    \lean{MPOTensor.PhysicalSectorFactorization.NeighboringTraceFactorization.completedNeighboringControlledMap_sameBlock_apply}
    \lean{MPOTensor.PhysicalSectorFactorization.NeighboringTraceFactorization.completedNeighboringPreparationMap_smul}
    \lean{MPOTensor.PhysicalSectorFactorization.NeighboringTraceFactorization.s1Map}
    \lean{MPOTensor.PhysicalSectorFactorization.NeighboringTraceFactorization.s1Map_block_eq_kronecker,
      MPOTensor.PhysicalSectorFactorization.NeighboringTraceFactorization.s1Map_apply_of_ne}
    \leanok
    \uses{def:mpdo_neighboring_trace_factorization}
    On a sector pair with $a_kb_h\ne0$, let
    \[
      \mathcal S_{k,h}(X)=X\otimes\frac{\eta_{k,h}}{a_kb_h}.
    \]
    If $a_kb_h=0$, positivity and vanishing trace imply
    $\eta_{k,h}=0$; choose any density operator on that unused summand.
    Taking the orthogonal direct sum over $(k,h)$ defines $\mathcal S_1$.
\end{definition}

\begin{theorem}[$\mathcal S_1$ is trace-preserving completely positive]
    \label{thm:mpdo_physical_sector_s1_map_cptp}
    \lean{MPOTensor.PhysicalSectorFactorization.NeighboringTraceFactorization.completedNeighboringControlledMap_isKrausCPTP}
    \lean{MPOTensor.PhysicalSectorFactorization.NeighboringTraceFactorization.s1Map_isKrausCPTP}
    \leanok
    \uses{def:mpdo_physical_sector_s1_map}
    The map $\mathcal S_1$ is trace-preserving and completely positive.
\end{theorem}

\begin{definition}[The coarse-graining channel $\mathcal S$]
    \label{def:mpdo_physical_sector_s_map}
    \lean{MPOTensor.PhysicalSectorFactorization.NeighboringTraceFactorization.sMap}
    \leanok
    \uses{def:mpdo_s0_map, def:mpdo_physical_sector_s1_map,
      def:mpdo_s2_map}
    Define the channel from three physical sites to two physical sites by
    \[
      \mathcal S=\mathcal S_2\mathcal S_1\mathcal S_0.
    \]
    This is the coarse-graining map in Proposition~C.7
    \cite[Appendix~C.2, lines~1547--1563]{Cirac2016MPDO_arXiv}.
\end{definition}

\begin{theorem}[$\mathcal S$ is trace-preserving completely positive]
    \label{thm:mpdo_physical_sector_s_map_cptp}
    \lean{MPOTensor.PhysicalSectorFactorization.NeighboringTraceFactorization.sMap_isKrausCPTP}
    \leanok
    \uses{thm:mpdo_s0_map_cptp,
      thm:mpdo_physical_sector_s1_map_cptp,
      thm:mpdo_s2_map_cptp,
      lem:is_kraus_cptp_comp}
    The channel $\mathcal S$ is trace-preserving and completely positive.
\end{theorem}

\begin{proof}
    \leanok
    \uses{lem:is_kraus_cptp_comp}
    Each of the three factors is trace-preserving and completely positive,
    and this property is preserved under composition.
\end{proof}

\begin{definition}[Two- and three-site sector coordinates]
    \label{def:mpdo_physical_sector_closure_coordinates}
    \lean{MPOTensor.PhysicalSectorFactorization.twoSiteSectorCoordinateEquiv,
      MPOTensor.PhysicalSectorFactorization.twoSiteSectorCoordinateEquiv_symm_apply,
      MPOTensor.PhysicalSectorFactorization.threeSiteSectorCoordinateEquiv,
      MPOTensor.PhysicalSectorFactorization.threeSiteSectorCoordinateEquiv_symm_apply}
    \leanok
    \uses{def:mpdo_two_site_sector_closure,
      def:mpdo_three_site_sector_closure}
    The coordinate equivalences used by the coarse-graining channel expose
    the direct sums over site-sector pairs and triples. Write
    $\mathfrak R_2$ and $\mathfrak R_3$ for the corresponding matrix
    reindexings. Their inverses send every sector block to the corresponding
    physical matrix entries. These reindexings preserve the left-right order
    within every physical site. They differ from the regrouped coordinates
    $\mathfrak G_2$ and $\mathfrak G_3$ of
    Definition~\ref{def:mpdo_global_sector_closure_coordinates}, which place
    the two outer subspins before the neighboring subspins. If
    $R_{\Phi_2}$ and $R_{\Phi_3}$ denote the two regrouping maps, then
    \[
      \mathfrak G_2=R_{\Phi_2}\mathfrak R_2,
      \qquad
      \mathfrak G_3=R_{\Phi_3}\mathfrak R_3.
    \]
\end{definition}

\begin{theorem}[Diagonal outer-sector block of the three-site closure]
    \label{thm:mpdo_regrouped_three_site_closure_block}
    \lean{MPOTensor.PhysicalSectorFactorization.s0Regrouped_threeSiteClosure_block_eq}
    \leanok
    \uses{def:phys_close,
      def:mpdo_sector_boundary_operator,
      def:mpdo_physical_sector_three_site_neighboring_operator,
      def:mpdo_physical_sector_closure_coordinates,
      def:mpdo_three_site_subspin_regrouping}
    For every virtual matrix $X$ and outer-sector pair $(k,h)$,
    \[
      \left[
        R_{\Phi_3}\mathfrak R_3({\cal K}_3(X))
      \right]_{(k,h),(k,h)}
      =B_{k,h}(X)\otimes\Omega_{k,h}.
    \]
\end{theorem}

\begin{proof}
    \leanok
    \uses{thm:mpdo_three_site_sector_closure_factorization,
      def:mpdo_physical_sector_factorization}
    On the summand labelled by the middle sector $l$, this is the fixed-sector
    factorization
    $B_{k,h}(X)\otimes(\eta_{k,l}\otimes\eta_{l,h})$.
    Taking the direct sum over $l$ gives the stated block.
\end{proof}

\begin{theorem}[Off-diagonal outer-sector entries of the three-site closure]
    \label{thm:mpdo_regrouped_three_site_closure_off_diagonal}
    \lean{MPOTensor.PhysicalSectorFactorization.s0Regrouped_threeSiteClosure_apply_of_outer_ne}
    \leanok
    \uses{def:mpdo_physical_sector_factorization,
      def:mpdo_physical_sector_closure_coordinates,
      def:mpdo_three_site_subspin_regrouping}
    Every matrix entry of
    $R_{\Phi_3}\mathfrak R_3({\cal K}_3(X))$ between distinct outer-sector
    pairs $(k,h)\ne(p,q)$ vanishes.
\end{theorem}

\begin{proof}
    \leanok
    \uses{def:mpdo_physical_sector_factorization}
    Distinct outer-sector pairs differ either in their first sector or in
    their third sector. In the corresponding physical slice, the
    sector-coordinate tensor is block diagonal, so that factor vanishes.
\end{proof}

\begin{theorem}[Regrouped two-site closure block]
    \label{thm:mpdo_t0_regrouped_two_site_closure_block}
    \lean{MPOTensor.PhysicalSectorFactorization.t0Regrouped_twoSiteClosure_block_eq}
    \leanok
    \uses{thm:mpdo_two_site_sector_closure_factorization,
      def:mpdo_physical_sector_closure_coordinates,
      def:mpdo_two_site_subspin_regrouping}
    For every virtual matrix $X$ and outer-sector pair $(k,h)$,
    \[
      \left[
        R_{\Phi_2}\mathfrak R_2({\cal K}_2(X))
      \right]_{(k,h),(k,h)}
      =B_{k,h}(X)\otimes\eta_{k,h}.
    \]
\end{theorem}

\begin{proof}
    \leanok
    \uses{thm:mpdo_two_site_sector_closure_factorization}
    This is the fixed-sector two-site factorization after identifying the
    sitewise sector coordinates with the regrouped boundary and neighboring
    coordinates.
\end{proof}

\begin{theorem}[Diagonal-block action of $\mathcal T_0$]
    \label{thm:mpdo_t0_same_block_action}
    \lean{MPOTensor.PhysicalSectorFactorization.t0Map_sameBlock_apply}
    \leanok
    \uses{def:mpdo_t0_map,
      def:mpdo_two_site_subspin_regrouping}
    For every two-site matrix $Y$, outer-sector pair $(k,h)$, and boundary
    indices $a,b$,
    \[
      [\mathcal T_0(Y)]_{(k,h,a),(k,h,b)}
      =\left[
        \tr_{R_k\otimes L_h}
          \left([R_{\Phi_2}Y]_{(k,h),(k,h)}\right)
       \right]_{a,b}.
    \]
\end{theorem}

\begin{proof}
    \leanok
    The controlled partial trace acts on a diagonal outer-sector block by
    the ordinary partial trace over its neighboring factor.
\end{proof}

\begin{theorem}[Partial trace of a factorized two-site block]
    \label{thm:mpdo_t0_factorized_block_action}
    \lean{MPOTensor.PhysicalSectorFactorization.NeighboringTraceFactorization.t0Map_block_eq_smul}
    \leanok
    \uses{thm:mpdo_t0_same_block_action,
      def:mpdo_neighboring_trace_factorization}
    If the $(k,h)$ diagonal block of $R_{\Phi_2}Y$ is
    $B\otimes\eta_{k,h}$, then
    \[
      [\mathcal T_0(Y)]_{(k,h),(k,h)}=a_kb_hB.
    \]
\end{theorem}

\begin{proof}
    \leanok
    \uses{thm:mpdo_t0_same_block_action,
      def:mpdo_neighboring_trace_factorization}
    The partial trace of $B\otimes\eta_{k,h}$ is
    $\tr(\eta_{k,h})B=a_kb_hB$.
\end{proof}

% **Local fix (zero-weight quotient):** the next three entries use the
% completion documented in
% docs/paper-gaps/cpgsv17_mpdo_zero_weight_preparation_completion.tex.
\begin{theorem}[Preparation of a scaled boundary block]
    \label{thm:mpdo_t1_scaled_block_preparation}
    \lean{MPOTensor.PhysicalSectorFactorization.NeighboringTraceFactorization.completedThreeSitePreparationMap_smul}
    \leanok
    \uses{def:mpdo_physical_sector_t1_preparation,
      thm:mpdo_physical_sector_t1_density_properties}
    For every outer-sector pair $(k,h)$ and boundary matrix $B$, the completed
    sectorwise preparation satisfies
    \[
      \overline{\mathcal T}_{k,h}(a_kb_hB)
      =B\otimes\Omega_{k,h}.
    \]
\end{theorem}

\begin{proof}
    \leanok
    \uses{thm:mpdo_physical_sector_t1_density_properties,
      thm:mpdo_physical_sector_three_site_neighboring_operator_trace}
    If $a_kb_h\ne0$, the scalar cancels the normalization in
    $\widehat\Omega_{k,h}=(a_kb_h)^{-1}\Omega_{k,h}$. If $a_kb_h=0$, both
    sides vanish because $\Omega_{k,h}=0$.
\end{proof}

\begin{theorem}[Diagonal-block action of $\mathcal T_1$]
    \label{thm:mpdo_t1_factorized_block_action}
    \lean{MPOTensor.PhysicalSectorFactorization.NeighboringTraceFactorization.t1Map_block_eq_kronecker}
    \leanok
    \uses{def:mpdo_physical_sector_t1_preparation,
      thm:mpdo_t1_scaled_block_preparation}
    If the $(k,h)$ diagonal block of a boundary matrix $Y$ is $a_kb_hB$, then
    \[
      [\mathcal T_1(Y)]_{(k,h),(k,h)}
      =B\otimes\Omega_{k,h}.
    \]
\end{theorem}

\begin{proof}
    \leanok
    \uses{thm:mpdo_t1_scaled_block_preparation}
    On a diagonal outer-sector block, the controlled map is the completed
    sectorwise preparation, to which the preceding theorem applies.
\end{proof}

\begin{theorem}[Off-diagonal action of $\mathcal T_1$]
    \label{thm:mpdo_t1_off_diagonal_action}
    \lean{MPOTensor.PhysicalSectorFactorization.NeighboringTraceFactorization.t1Map_apply_of_ne}
    \leanok
    \uses{def:mpdo_physical_sector_t1_preparation}
    The map $\mathcal T_1$ annihilates every matrix entry between distinct
    outer-sector pairs.
\end{theorem}

\begin{proof}
    \leanok
    This is the orthogonal control in the definition of $\mathcal T_1$.
\end{proof}

\begin{theorem}[Reindexing action of $\mathcal T_2$]
    \label{thm:mpdo_t2_reindexing_action}
    \lean{MPOTensor.PhysicalSectorFactorization.t2Map_apply}
    \leanok
    \uses{def:mpdo_t2_map}
    Let $\Phi_3$ be the three-site regrouping. For every prepared matrix $Y$
    and three-site indices $p,q$,
    \[
      [\mathcal T_2(Y)]_{p,q}
      =Y_{\Phi_3(p),\Phi_3(q)}.
    \]
\end{theorem}

\begin{proof}
    \leanok
    This is matrix reindexing along the inverse regrouping
    $\Phi_3^{-1}$.
\end{proof}

\begin{theorem}[Refinement of the two-site closure]
    \label{thm:mpdo_physical_sector_t_closure_identity}
    \lean{MPOTensor.PhysicalSectorFactorization.NeighboringTraceFactorization.tMap_twoSiteClosure_eq_threeSiteClosure}
    \leanok
    \uses{def:phys_close,
      def:mpdo_physical_sector_t_map,
      def:mpdo_physical_sector_closure_coordinates,
      def:mpdo_neighboring_trace_factorization}
    Assume the neighboring-operator trace factorization of
    Definition~\ref{def:mpdo_neighboring_trace_factorization}. For every
    virtual matrix $X$, in the canonical physical-sector coordinates selected
    by the isometry $U$, the refinement channel satisfies
    \[
      \mathcal T\bigl(\mathfrak R_2({\cal K}_2(X))\bigr)
      =\mathfrak R_3({\cal K}_3(X)).
    \]
    This is the refinement half of Proposition~C.7
    \cite[Appendix~C.2, lines~1510--1516 and~1522--1545]
      {Cirac2016MPDO_arXiv}.
\end{theorem}

\begin{proof}
    \leanok
    \uses{thm:mpdo_t0_regrouped_two_site_closure_block,
      thm:mpdo_t0_factorized_block_action,
      thm:mpdo_t1_factorized_block_action,
      thm:mpdo_t1_off_diagonal_action,
      thm:mpdo_regrouped_three_site_closure_block,
      thm:mpdo_regrouped_three_site_closure_off_diagonal}
    The regrouping contained in $\mathcal T_0$ changes
    $\mathfrak R_2$ into $\mathfrak G_2$. On the $(k,h)$ diagonal block, the
    partial trace then gives
    \[
      \mathcal T_0\bigl(\mathfrak R_2({\cal K}_2(X))\bigr)_{k,h}
      =a_kb_hB_{k,h}(X).
    \]
    On an active pair, $\mathcal T_1$ adjoins
    $(a_kb_h)^{-1}\Omega_{k,h}$, and hence
    \[
      \mathcal T_1\mathcal T_0
        \bigl(\mathfrak R_2({\cal K}_2(X))\bigr)_{k,h}
      =B_{k,h}(X)\otimes\Omega_{k,h}.
    \]
    If $a_kb_h=0$, both sides vanish because
    $\Omega_{k,h}=0$. After the canonical reassociation of the intermediate
    sector sum, these are precisely the diagonal blocks of
    $\mathfrak G_3({\cal K}_3(X))$; all off-diagonal sector blocks vanish.
    Finally, the inverse regrouping $\mathcal T_2$ changes
    $\mathfrak G_3$ into $\mathfrak R_3$, which gives the stated equality.
\end{proof}

\begin{theorem}[Coarse-graining of the three-site closure]
    \label{thm:mpdo_physical_sector_s_closure_identity}
    \lean{MPOTensor.PhysicalSectorFactorization.NeighboringTraceFactorization.s0Map_block_eq_smul}
    \lean{MPOTensor.PhysicalSectorFactorization.NeighboringTraceFactorization.sMap_threeSiteClosure_eq_twoSiteClosure}
    \leanok
    \uses{def:phys_close, def:mpdo_physical_sector_s_map,
      def:mpdo_physical_sector_closure_coordinates,
      def:mpdo_neighboring_trace_factorization}
    Assume the neighboring-operator trace factorization of
    Definition~\ref{def:mpdo_neighboring_trace_factorization}. For every
    virtual matrix $X$, in the canonical physical-sector
    coordinates selected by the isometry $U$, the coarse-graining channel
    satisfies
    \[
      \mathcal S\bigl(\mathfrak R_3({\cal K}_3(X))\bigr)
      =\mathfrak R_2({\cal K}_2(X)).
    \]
\end{theorem}

\begin{proof}
    \leanok
    \uses{thm:mpdo_regrouped_three_site_closure_block,
      thm:mpdo_two_site_sector_closure_factorization,
      thm:mpdo_three_site_sector_closure_factorization,
      thm:mpdo_physical_sector_three_site_neighboring_operator_trace,
      def:mpdo_physical_sector_factorization,
      def:mpdo_physical_sector_closure_coordinates,
      def:mpdo_physical_sector_s1_map}
    On the $(k,h)$ diagonal block, $\mathcal S_0$ traces
    $\Omega_{k,h}$ and gives $a_kb_hB_{k,h}(X)$. On an active sector pair,
    \[
      \mathcal S_1\bigl((a_kb_h)B_{k,h}(X)\bigr)
      =B_{k,h}(X)\otimes\eta_{k,h}.
    \]
    If $a_kb_h=0$, positivity and zero trace give $\eta_{k,h}=0$, so both
    sides of this identity vanish. Finally, $\mathcal S_2$ regroups the four
    factors into two physical sites. For distinct sector pairs
    $(k,h)\ne(p,q)$,
    \[
      (\mathcal S_1Y)_{(k,h),(p,q)}=0,
      \qquad
      [\mathfrak R_2({\cal K}_2(X))]_{(k,h),(p,q)}=0.
    \]
\end{proof}

For the following construction, through
Theorem~\ref{thm:mpdo_blocked_original_tensor_is_rfp_via_ts}, assume the
neighboring-operator trace factorization of
Definition~\ref{def:mpdo_neighboring_trace_factorization}. Write
\[
  q=\dim\!\left(\bigoplus_k L_k\otimes R_k\right)
\]
for the physical dimension of the sector-coordinate tensor
$\widehat{\cal K}$.

\begin{definition}[Refinement in ordinary physical coordinates]
    \label{def:mpdo_physical_t_map}
    \lean{MPOTensor.PhysicalSectorFactorization.NeighboringTraceFactorization.physicalTMap}
    \leanok
    \uses{def:mpdo_physical_sector_t_map,
      def:mpdo_physical_sector_closure_coordinates}
    Transporting $\mathcal T$ through the canonical two-site and three-site
    coordinate identifications gives a channel
    \[
      \widehat{\mathcal T}:M_{q^2}(\mathbb C)\longrightarrow
        M_{q^3}(\mathbb C),
    \]
    where all multiple physical indices are associated to the right.
\end{definition}

\begin{theorem}[The physical refinement map is a channel]
    \label{thm:mpdo_physical_t_map_cptp}
    \lean{MPOTensor.PhysicalSectorFactorization.NeighboringTraceFactorization.physicalTMap_isKrausCPTP}
    \leanok
    \uses{def:mpdo_physical_t_map}
    The map $\widehat{\mathcal T}$ is trace-preserving and completely
    positive.
\end{theorem}

\begin{proof}\leanok
    \uses{thm:mpdo_physical_sector_t_map_cptp}
    The coordinate identifications are unitary permutation channels.
    Composing them with a trace-preserving completely positive map preserves
    both properties.
\end{proof}

% **Local fix (zero-weight quotient):** this identity uses the completed
% zero-weight preparation. Documented in
% docs/paper-gaps/cpgsv17_mpdo_zero_weight_preparation_completion.tex.
\begin{theorem}[Refinement identity in ordinary physical coordinates]
    \label{thm:mpdo_physical_t_closure_identity}
    \lean{MPOTensor.PhysicalSectorFactorization.NeighboringTraceFactorization.physicalTMap_twoSiteClosure_eq_threeSiteClosure}
    \leanok
    \uses{def:mpdo_physical_t_map}
    For every virtual matrix $X$,
    \[
      \widehat{\mathcal T}\bigl(\widehat{\cal K}_2(X)\bigr)
      =\widehat{\cal K}_3(X).
    \]
\end{theorem}

\begin{proof}\leanok
    \uses{thm:mpdo_physical_sector_t_closure_identity}
    Transport the sector-coordinate identity through the inverse three-site
    coordinate identification.
\end{proof}

\begin{definition}[Physical coarse-graining]
    \label{def:mpdo_physical_s_map}
    \lean{MPOTensor.PhysicalSectorFactorization.NeighboringTraceFactorization.physicalSMap}
    \leanok
    \uses{def:mpdo_physical_sector_s_map,
      def:mpdo_physical_sector_closure_coordinates}
    Transporting $\mathcal S$ through the canonical three-site and two-site
    coordinate identifications gives a channel
    \[
      \widehat{\mathcal S}:M_{q^3}(\mathbb C)\longrightarrow
        M_{q^2}(\mathbb C).
    \]
\end{definition}

\begin{theorem}[The physical coarse-graining map is a channel]
    \label{thm:mpdo_physical_s_map_cptp}
    \lean{MPOTensor.PhysicalSectorFactorization.NeighboringTraceFactorization.physicalSMap_isKrausCPTP}
    \leanok
    \uses{def:mpdo_physical_s_map}
    The map $\widehat{\mathcal S}$ is trace-preserving and completely
    positive.
\end{theorem}

\begin{proof}\leanok
    \uses{thm:mpdo_physical_sector_s_map_cptp}
    This follows by composition with the two coordinate-permutation
    channels.
\end{proof}

\begin{theorem}[Coarse-graining identity in ordinary physical coordinates]
    \label{thm:mpdo_physical_s_closure_identity}
    \lean{MPOTensor.PhysicalSectorFactorization.NeighboringTraceFactorization.physicalSMap_threeSiteClosure_eq_twoSiteClosure}
    \leanok
    \uses{def:mpdo_physical_s_map}
    For every virtual matrix $X$,
    \[
      \widehat{\mathcal S}\bigl(\widehat{\cal K}_3(X)\bigr)
      =\widehat{\cal K}_2(X).
    \]
\end{theorem}

\begin{proof}\leanok
    \uses{thm:mpdo_physical_sector_s_closure_identity}
    Transport the sector-coordinate identity through the inverse two-site
    coordinate identification.
\end{proof}

\begin{definition}[Localized physical refinement]
    \label{def:mpdo_localized_physical_t_map}
    \lean{MPOTensor.PhysicalSectorFactorization.NeighboringTraceFactorization.localizedPhysicalTMap}
    \leanok
    \uses{def:mpdo_physical_t_map}
    Reassociating three right-associated sites as $(12)3$, applying
    $\widehat{\mathcal T}\otimes\id$ to the first two, and
    reassociating the result back to $1(2(34))$ defines the localized
    refinement channel
    \[
      \widehat{\mathcal T}_{(12)3}:
      M_{q^3}(\mathbb C)\longrightarrow M_{q^4}(\mathbb C)
    \]
    on right-associated physical coordinates.
\end{definition}

\begin{theorem}[The localized physical refinement is a channel]
    \label{thm:mpdo_localized_physical_t_map_cptp}
    \lean{MPOTensor.PhysicalSectorFactorization.NeighboringTraceFactorization.localizedPhysicalTMap_isKrausCPTP}
    \leanok
    \uses{def:mpdo_localized_physical_t_map}
    The localized refinement map is trace-preserving and completely
    positive.
\end{theorem}

\begin{proof}\leanok
    \uses{thm:mpdo_physical_t_map_cptp}
    Tensoring a channel with the identity preserves complete positivity and
    trace, as do the two associativity permutations.
\end{proof}

% **Local fix (zero-weight quotient):** this identity uses the completed
% zero-weight preparation. Documented in
% docs/paper-gaps/cpgsv17_mpdo_zero_weight_preparation_completion.tex.
\begin{theorem}[Localized refinement of physical closures]
    \label{thm:mpdo_localized_physical_t_closure_identity}
    \lean{MPOTensor.PhysicalSectorFactorization.NeighboringTraceFactorization.localizedPhysicalTMap_threeSiteClosure_eq_fourSiteClosure}
    \leanok
    \uses{def:mpdo_localized_physical_t_map}
    For every virtual matrix $X$,
    \[
      (\widehat{\mathcal T}\otimes\id)
        \bigl(\widehat{\cal K}_3(X)\bigr)
      =\widehat{\cal K}_4(X),
    \]
    with the displayed map understood in right-associated coordinates.
\end{theorem}

\begin{proof}\leanok
    \uses{thm:mpdo_physical_t_closure_identity}
    Fixing the last ket and bra indices turns the three-site closure into a
    two-site closure with the last tensor matrix absorbed into $X$. Apply the
    two-to-three-site refinement identity to this virtual matrix.
\end{proof}

\begin{definition}[Localized physical coarse-graining]
    \label{def:mpdo_localized_physical_s_map}
    \lean{MPOTensor.PhysicalSectorFactorization.NeighboringTraceFactorization.localizedPhysicalSMap}
    \leanok
    \uses{def:mpdo_physical_s_map}
    Reassociating four right-associated sites as $(123)4$, applying
    $\widehat{\mathcal S}\otimes\id$ to the first three, and
    reassociating the result back to $1(23)$ defines the localized
    coarse-graining channel
    \[
      \widehat{\mathcal S}_{(123)4}:
      M_{q^4}(\mathbb C)\longrightarrow M_{q^3}(\mathbb C)
    \]
    on right-associated physical coordinates.
\end{definition}

\begin{theorem}[The localized physical coarse-graining is a channel]
    \label{thm:mpdo_localized_physical_s_map_cptp}
    \lean{MPOTensor.PhysicalSectorFactorization.NeighboringTraceFactorization.localizedPhysicalSMap_isKrausCPTP}
    \leanok
    \uses{def:mpdo_localized_physical_s_map}
    The localized coarse-graining map is trace-preserving and completely
    positive.
\end{theorem}

\begin{proof}\leanok
    \uses{thm:mpdo_physical_s_map_cptp}
    Tensor the physical coarse-graining channel with the identity and compose
    with the two associativity permutations.
\end{proof}

\begin{theorem}[Localized coarse-graining of physical closures]
    \label{thm:mpdo_localized_physical_s_closure_identity}
    \lean{MPOTensor.PhysicalSectorFactorization.NeighboringTraceFactorization.localizedPhysicalSMap_fourSiteClosure_eq_threeSiteClosure}
    \leanok
    \uses{def:mpdo_localized_physical_s_map}
    For every virtual matrix $X$,
    \[
      (\widehat{\mathcal S}\otimes\id)
        \bigl(\widehat{\cal K}_4(X)\bigr)
      =\widehat{\cal K}_3(X),
    \]
    in right-associated coordinates.
\end{theorem}

\begin{proof}\leanok
    \uses{thm:mpdo_physical_s_closure_identity}
    Fixing the last ket and bra indices identifies the four-site closure with
    a three-site closure having the last tensor matrix absorbed into $X$.
    Apply the three-to-two-site coarse-graining identity.
\end{proof}

% **Local fix (implication label):** Appendix line 1824 says
% (iii) implies (v), but the theorem statement and the local structure used
% here give (iv) implies (v). Documented in
% docs/paper-gaps/cpgsv17_mpdo_theorem_4_9_implication_label.tex.
\begin{definition}[Two-step physical refinement]
    \label{def:mpdo_physical_t_squared_map}
    \lean{MPOTensor.PhysicalSectorFactorization.NeighboringTraceFactorization.physicalTSquaredMap}
    \leanok
    \uses{def:mpdo_physical_t_map,
      def:mpdo_localized_physical_t_map}
    Let $\widehat{\cal K}$ be the sector-coordinate tensor selected by the
    physical isometry of Proposition~C.7. Define
    \[
      \widehat{\mathcal T}^2
      =\widehat{\mathcal T}_{(12)3}\circ\widehat{\mathcal T},
      \qquad
      \widehat{\mathcal T}^2:M_{q^2}(\mathbb C)
        \longrightarrow M_{q^4}(\mathbb C),
    \]
    where the second application acts on the first two sites and leaves the
    third site unchanged. Thus the superscript denotes two successive
    overlapping applications, as in the proof of
    \cite[Theorem~4.9, (iv)$\Rightarrow$(v)]{Cirac2016MPDO_arXiv}.
\end{definition}

\begin{theorem}[The two-step physical refinement is a channel]
    \label{thm:mpdo_physical_t_squared_map_cptp}
    \lean{MPOTensor.PhysicalSectorFactorization.NeighboringTraceFactorization.physicalTSquaredMap_isKrausCPTP}
    \leanok
    \uses{def:mpdo_physical_t_squared_map}
    The map $\widehat{\mathcal T}^2$ is trace-preserving and completely
    positive.
\end{theorem}

\begin{proof}\leanok
    \uses{thm:mpdo_physical_t_map_cptp,
      thm:mpdo_localized_physical_t_map_cptp}
    It is the composition of the physical refinement channel and its
    localization on the first two sites.
\end{proof}

% **Local fix (zero-weight quotient):** this identity uses the completed
% zero-weight preparation. Documented in
% docs/paper-gaps/cpgsv17_mpdo_zero_weight_preparation_completion.tex.
\begin{theorem}[Two-step refinement of the sector-coordinate tensor]
    \label{thm:mpdo_physical_t_squared_closure_identity}
    \lean{MPOTensor.PhysicalSectorFactorization.NeighboringTraceFactorization.physicalTSquaredMap_twoSiteClosure_eq_fourSiteClosure}
    \leanok
    \uses{def:mpdo_physical_t_squared_map}
    For every virtual matrix $X$, the sector-coordinate tensor satisfies
    \[
      \widehat{\mathcal T}^2
        \bigl(\widehat{\cal K}_2(X)\bigr)
      =\widehat{\cal K}_4(X).
    \]
\end{theorem}

\begin{proof}\leanok
    \uses{thm:mpdo_physical_t_closure_identity,
      thm:mpdo_localized_physical_t_closure_identity}
    The first refinement gives $\widehat{\cal K}_3(X)$; applying the localized
    refinement to its first two sites gives $\widehat{\cal K}_4(X)$.
\end{proof}

\begin{definition}[Two-step physical coarse-graining]
    \label{def:mpdo_physical_s_squared_map}
    \lean{MPOTensor.PhysicalSectorFactorization.NeighboringTraceFactorization.physicalSSquaredMap}
    \leanok
    \uses{def:mpdo_physical_s_map,
      def:mpdo_localized_physical_s_map}
    Define
    \[
      \widehat{\mathcal S}^2
      =\widehat{\mathcal S}\circ\widehat{\mathcal S}_{(123)4},
      \qquad
      \widehat{\mathcal S}^2:M_{q^4}(\mathbb C)
        \longrightarrow M_{q^2}(\mathbb C),
    \]
    where the first application acts on the first three sites and leaves the
    fourth site unchanged. The superscript again denotes successive
    overlapping applications.
\end{definition}

\begin{theorem}[The two-step physical coarse-graining is a channel]
    \label{thm:mpdo_physical_s_squared_map_cptp}
    \lean{MPOTensor.PhysicalSectorFactorization.NeighboringTraceFactorization.physicalSSquaredMap_isKrausCPTP}
    \leanok
    \uses{def:mpdo_physical_s_squared_map}
    The map $\widehat{\mathcal S}^2$ is trace-preserving and completely
    positive.
\end{theorem}

\begin{proof}\leanok
    \uses{thm:mpdo_physical_s_map_cptp,
      thm:mpdo_localized_physical_s_map_cptp}
    It is the composition of the localized coarse-graining channel and the
    physical coarse-graining channel.
\end{proof}

\begin{theorem}[Two-step coarse-graining of the sector-coordinate tensor]
    \label{thm:mpdo_physical_s_squared_closure_identity}
    \lean{MPOTensor.PhysicalSectorFactorization.NeighboringTraceFactorization.physicalSSquaredMap_fourSiteClosure_eq_twoSiteClosure}
    \leanok
    \uses{def:mpdo_physical_s_squared_map}
    For every virtual matrix $X$, the sector-coordinate tensor satisfies
    \[
      \widehat{\mathcal S}^2
        \bigl(\widehat{\cal K}_4(X)\bigr)
      =\widehat{\cal K}_2(X).
    \]
\end{theorem}

\begin{proof}\leanok
    \uses{thm:mpdo_localized_physical_s_closure_identity,
      thm:mpdo_physical_s_closure_identity}
    The localized coarse-graining first gives $\widehat{\cal K}_3(X)$; the
    second coarse-graining gives $\widehat{\cal K}_2(X)$.
\end{proof}

\begin{figure}[!b]
    \centering
    % Formula: \widehat{\mathcal T}(\widehat{\cal K}_2(X))
    % =\widehat{\cal K}_3(X), followed by
    % \widehat{\mathcal T}_{(12)3}(\widehat{\cal K}_3(X))
    % =\widehat{\cal K}_4(X); the reverse row is the corresponding pair of
    % \widehat{\mathcal S} identities.
    % Ink-to-index: every \widehat{\cal K} box has west/east virtual indices
    % and upper/lower physical indices.  Each dashed span groups two adjacent
    % tensor bodies into one blocked local site and introduces no contraction.
    % Contractions: only adjacent east--west virtual legs are summed.  All
    % upper and lower physical legs, and both outer virtual legs, remain free.
    % Boundary signatures (virtual-west, virtual-east, physical-up,
    % physical-down): the two-, three-, and four-site panels have respectively
    % (1,1,2,2), (1,1,3,3), and (1,1,4,4).  The arrows are physical channels,
    % so their endpoint physical arities change as displayed.
    % Source verdict: Appendix C.2, lines 1821--1825, specifies the two-step
    % channels algebraically but has no dedicated figure.  The rows reproduce
    % that composition using the local tensors already defined here.
    \begin{tenkzequation}
        $\widehat{\mathcal T}^{\,2}:\quad$
        \tnpic[compact, physical=updown, tensor style=box]{%
            \tn[mpo]{\widehat{\cal K}}
                \tnspan[box, label pos=south]{2}{\widehat{\cal K}^{[2]}} &
            \tn[mpo]{\widehat{\cal K}}}
        $\quad\xrightarrow{\widehat{\mathcal T}}\quad$
        \tnpic[compact, physical=updown, tensor style=box]{%
            \tn[mpo]{\widehat{\cal K}} &
            \tn[mpo]{\widehat{\cal K}} &
            \tn[mpo]{\widehat{\cal K}}}
        $\quad\xrightarrow{\widehat{\mathcal T}_{(12)3}}\quad$
        \tnpic[compact, physical=updown, tensor style=box]{%
            \tn[mpo]{\widehat{\cal K}}
                \tnspan[box, label pos=south]{2}{\widehat{\cal K}^{[2]}} &
            \tn[mpo]{\widehat{\cal K}} &
            \tn[mpo]{\widehat{\cal K}}
                \tnspan[box, label pos=south]{2}{\widehat{\cal K}^{[2]}} &
            \tn[mpo]{\widehat{\cal K}}}
    \end{tenkzequation}
    \par\medskip
    \begin{tenkzequation}
        $\widehat{\mathcal S}^{\,2}:\quad$
        \tnpic[compact, physical=updown, tensor style=box]{%
            \tn[mpo]{\widehat{\cal K}}
                \tnspan[box, label pos=south]{2}{\widehat{\cal K}^{[2]}} &
            \tn[mpo]{\widehat{\cal K}} &
            \tn[mpo]{\widehat{\cal K}}
                \tnspan[box, label pos=south]{2}{\widehat{\cal K}^{[2]}} &
            \tn[mpo]{\widehat{\cal K}}}
        $\quad\xrightarrow{\widehat{\mathcal S}_{(123)4}}\quad$
        \tnpic[compact, physical=updown, tensor style=box]{%
            \tn[mpo]{\widehat{\cal K}} &
            \tn[mpo]{\widehat{\cal K}} &
            \tn[mpo]{\widehat{\cal K}}}
        $\quad\xrightarrow{\widehat{\mathcal S}}\quad$
        \tnpic[compact, physical=updown, tensor style=box]{%
            \tn[mpo]{\widehat{\cal K}}
                \tnspan[box, label pos=south]{2}{\widehat{\cal K}^{[2]}} &
            \tn[mpo]{\widehat{\cal K}}}
    \end{tenkzequation}
    \caption{The superscripts in $\widehat{\mathcal T}^2$ and
    $\widehat{\mathcal S}^2$ denote successive overlapping applications.
    The dashed frames identify the two-site and four-site endpoints with one
    and two sites of the blocked tensor. This is the channel construction in
    \cite[Appendix~C.2, lines~1821--1825]{Cirac2016MPDO_arXiv}.}
    \label{fig:mpdo_blocked_rfp_channels}
\end{figure}

% **Local fix (zero-weight quotient):** this theorem uses the completed
% zero-weight preparation through the two-step refinement identity. Documented
% in docs/paper-gaps/cpgsv17_mpdo_zero_weight_preparation_completion.tex.
% **Scope restriction (physical coordinates):** this theorem concerns the
% sector-coordinate tensor, not yet the original tensor. Documented in
% docs/paper-gaps/cpgsv17_mpdo_blocked_rfp_physical_transport.tex.
\begin{theorem}[Blocked sector-coordinate tensor is a renormalization fixed point]
    \label{thm:mpdo_blocked_sector_coordinate_tensor_is_rfp_via_ts}
    \lean{MPOTensor.PhysicalSectorFactorization.NeighboringTraceFactorization.blockTwo_sectorCoordinateTensor_isRFPViaTS}
    \leanok
    \uses{def:is_rfp_via_ts, def:mpdo_two_site_blocking}
    Let $\widehat{\cal K}^{[2]}$ be the tensor obtained by blocking two
    adjacent sites of the sector-coordinate tensor $\widehat{\cal K}$. Then
    $\widehat{\cal K}^{[2]}$ is a renormalization fixed point in the sense of
    Definition~\ref{def:is_rfp_via_ts}. More precisely, after the canonical
    identifications of one blocked site with two original sites and two
    blocked sites with four original sites, the required channels are
    $\widehat{\mathcal S}^2$ and
    $\widehat{\mathcal T}^2$.
\end{theorem}

\begin{proof}\leanok
    \uses{thm:mpdo_physical_t_squared_map_cptp,
      thm:mpdo_physical_t_squared_closure_identity,
      thm:mpdo_physical_s_squared_map_cptp,
      thm:mpdo_physical_s_squared_closure_identity,
      thm:mpdo_one_site_closure_two_site_blocking,
      thm:mpdo_two_site_closure_two_site_blocking}
    Decode one blocked index as a pair of physical indices and two blocked
    indices as four physical indices. The one-site and two-site closures of
    $\widehat{\cal K}^{[2]}$ thereby become respectively
    $\widehat{\cal K}_2(X)$ and $\widehat{\cal K}_4(X)$. Transporting the two
    channels through these permutation identifications preserves complete
    positivity and trace, and the two closure identities give the required
    equations in Definition~\ref{def:is_rfp_via_ts}.
\end{proof}

% **Local fix (zero-weight quotient):** this theorem uses the completed
% zero-weight preparation through the two-step refinement identity. Documented
% in docs/paper-gaps/cpgsv17_mpdo_zero_weight_preparation_completion.tex.
% **Local fix (implication label):** Appendix line 1824 says
% (iii) implies (v), but the theorem statement and the local structure used
% here give (iv) implies (v). Documented in
% docs/paper-gaps/cpgsv17_mpdo_theorem_4_9_implication_label.tex.
% **Scope restriction (SAL and ZCL):** the theorem starts from a neighboring
% trace factorization, including the rank-one identity
% tr(eta_{k,h}) = a_k b_h. It proves the conditional implication (iv) => (v),
% not Proposition prop2to5 from SAL and ZCL alone. The missing rank-one step is
% documented in docs/paper-gaps/cpgsv17_pf_rank_one.tex.
\begin{theorem}[Two-site blocking of the original tensor is a renormalization fixed point]
    \label{thm:mpdo_blocked_original_tensor_is_rfp_via_ts}
    \lean{MPOTensor.PhysicalSectorFactorization.NeighboringTraceFactorization.blockTwo_isRFPViaTS}
    \leanok
    \uses{def:is_rfp_via_ts, def:mpdo_two_site_blocking}
    Let ${\cal K}^{[2]}$ be the tensor obtained by blocking two adjacent sites
    of the original tensor ${\cal K}$. Then ${\cal K}^{[2]}$ is a
    renormalization fixed point in the sense of
    Definition~\ref{def:is_rfp_via_ts}. This is
    \cite[Theorem~4.9, (iv)$\Rightarrow$(v)]{Cirac2016MPDO_arXiv}.
\end{theorem}

\begin{proof}\leanok
    \uses{thm:mpdo_blocked_sector_coordinate_tensor_is_rfp_via_ts}
    Transport the refinement and coarse-graining channels for
    $\widehat{\cal K}^{[2]}$ through the two-site and four-site tensor powers
    of the physical isometry, as in
    equation~(\ref{eq:mpdo_physical_isometry_transport}). These unitary
    congruences preserve complete positivity and trace, while the
    corresponding closure identities carry the two fixed-point equations back
    to ${\cal K}^{[2]}$.
\end{proof}

% **Refuted prerequisite (Lemma C.5):**
% lem:mpdo_sal_zcl_rank_one_trace_matrix is formally refuted under its printed
% SAL and ZCL hypotheses by thm:mpdo_sal_zcl_rank_one_counterexample. Hence the
% printed proof through prop2to3 does not establish implication (ii) => (v),
% which remains unformalized rather than refuted. A conditional construction of
% rank-one and local-structure data under positive semidefiniteness or linear
% independence is given by thm:simple_mpdo_local_structure_data_of_sal_zcl. The
% conditional blocking theorem thm:mpdo_blocked_original_tensor_is_rfp_via_ts
% instead assumes a neighboring trace factorization.
% Documented in docs/paper-gaps/cpgsv17_pf_rank_one.tex.
\begin{theorem}[SAL and ZCL give the two blocking channels]
    \label{thm:mpdo_sal_zcl_blocking_channels}
    \notready
    \uses{def:is_sal, def:mpo_source_zcl, def:is_rfp_via_ts,
        def:mpdo_two_site_blocking,
        thm:mpdo_sal_zcl_bnt_sector_structure}
    Let $K$ be a simple tensor in block-injective canonical form (biCF),
    equipped with a basis-of-normal-tensors decomposition and satisfying the
    strong area law and zero correlation length. Let $M=K^{[2]}$ be its
    two-site blocking. Then
    there are trace-preserving completely positive maps in both directions
    between the one-site and two-site closures of $M$. Equivalently, $M$ is a
    renormalization fixed point via these two maps. This is implication
    (ii)$\Rightarrow$(v) of
    \cite[Theorem~4.9 and Appendix~C.2, lines~1810--1825]
      {Cirac2016MPDO_arXiv}. The separate standing biCF and BNT hypotheses are
    imposed at lines~849--850 and restated at line~1628 of the same source.
\end{theorem}

% Correct equation, not a figure: the cited source draws the two-to-three
% channel MPDO_TSKKK.png after explicitly suppressing U, but it does not draw
% this tensor-power transport square.  The transport theorem is recorded in
% docs/paper-gaps/cpgsv17_mpdo_blocked_rfp_physical_transport.tex.
\begin{equation}
    \label{eq:mpdo_physical_isometry_transport}
        U^{\otimes4}\widehat{\mathcal T}^{\,2}(U^\dagger)^{\otimes2}=\mathcal T,\qquad
        U^{\otimes2}\widehat{\mathcal S}^{\,2}(U^\dagger)^{\otimes4}=\mathcal S
\end{equation}
Conjugation by $U^{\otimes n}$ and $(U^\dagger)^{\otimes n}$ carries the
sector-coordinate channels to the ordinary-coordinate channels. This combines
Proposition~C.7 with
\cite[Appendix~C.2, lines~1510--1563 and 1821--1825]{Cirac2016MPDO_arXiv}.

Tracing the ket and bra indices of a sector tensor gives the bond vector
$|l_k)$ or the bond functional $(r_k|$.  Hence a neighboring pair gives
\[
    T_{k,h}=(r_k|l_h)=\tr(\eta_{k,h}),
\]
as in \cite[Appendix~C.2, lines~1474--1481]{Cirac2016MPDO_arXiv}.
